\section{Parameter selection for the sharp upper bound}\label{app:parameters}

This appendix records the order of choices for the proof of \cref{thm:sharp-upper}. Fix $\rho>0$ and the interior analytic radius $\rho/2$ used in \cref{lem:moment-transfer}. Let $\omega_0=\omega_{1/2}>0$ be the harmonic-measure bound from \cref{lem:two-constants}; the same value is used in both propagation steps. Choose $A>8$ and then choose $c_0=c_0(A,\rho)>0$ sufficiently small that the Cauchy estimate in \cref{lem:moment-transfer} dominates all factorial losses; equivalently, after absorbing the fixed $\rho$-dependent constants, one may require
\[
(A+3)c_0<\omega_0.
\]
With $m=\lfloor c_0L/\log L\rfloor$, this gives the required super-exponential moment accuracy.

Fix a finite target exponent $\Gamma>4$. Choose the truncation factor $B$ sufficiently large that the tail-lifting estimate has exponent $\beta_B>4\Gamma$. After $B$ is fixed, let $C_T$ denote the constant in the real-axis Taylor remainder and choose $a_0>0$ so small that
\[
C_Ta_0^2<e^{-4\Gamma},\qquad
4(1-\omega_0)(Ba_0+a_0^2)<\frac{\omega_0\Gamma}{2},\qquad
a_0^2<\frac{\omega_0\Gamma}{4}.
\]
Each Taylor remainder is at most $e^{-4\Gamma m}$, and the truncation contribution is at most $4Ce^{-\beta_Bm}$. Together with the super-exponential moment mismatch, their sum is at most $e^{-\Gamma m}$ for large $m$. Thus \cref{lem:real-window} uses the target exponent $\Gamma$, while the complex propagation estimate has raw exponent
\[
\gamma_{\mathrm{raw}}
:=\omega_0\Gamma-4(1-\omega_0)(Ba_0+a_0^2)
\ge\frac{\omega_0\Gamma}{2}>a_0^2.
\]
Choose any fixed $\gamma_*$ with
\[
a_0^2<\gamma_*<\gamma_{\mathrm{raw}}.
\]
The $O(1)$ term in the two-constants estimate is then absorbed for all sufficiently large $m$, giving the exponent used in \cref{lem:zero-free}. Finally choose the fixed Esseen fraction $0<\theta<1/4$ so that the cubic remainder in the logarithmic expansion is at most one quarter of the Gaussian quadratic damping on $|t|\le\theta R$.

The choices are made in the order
\[
A\to c_0\to\Gamma\to B\to a_0\to\gamma_*\to\theta,
\]
and no later choice affects an earlier inequality.
